\chapter{State-Space Methods and Modern Control}
\label{chap:modern}

\begin{learningobjective}
\begin{itemize}[leftmargin=*]
    \item Evaluate controllability and observability rigorously and understand their physical meaning.
    \item Design state-feedback controllers via pole placement and Linear Quadratic Regulation (LQR), including servo prefilters.
    \item Integrate state estimation (Kalman filtering) to close the loop in the presence of measurement noise.
\end{itemize}
\end{learningobjective}

\section{Concept Map for Beginners}
\begin{conceptroadmap}
\begin{enumerate}[leftmargin=*]
    \item Treat the state as the story of the system between input and output; once that narrative is clear, matrices simply formalize it.
    \item Diagnose controllability and observability before designing anything so you know which goals are achievable.
    \item Pair optimal control (LQR) with optimal estimation (Kalman filter) to build the modern separation-principle controller seen in practice.
\end{enumerate}
\end{conceptroadmap}

\section{Concept--Math--Engineering Loop}
\begin{table}[h]
    \centering
    \small
    \renewcommand{\arraystretch}{1.2}
    \begin{tabularx}{\linewidth}{|>{\raggedright\arraybackslash}p{0.28\linewidth}|>{\raggedright\arraybackslash}p{0.34\linewidth}|>{\raggedright\arraybackslash}X|}
        \hline
        \textbf{Concept} & \textbf{Mathematics} & \textbf{Engineering Practice} \\
        \hline
        Keeping the state explicit & Canonical realisations, matrix assembly & Recreate $A$, $B$, $C$, $D$ from a transfer function with \texttt{control.ss} to confirm intuition. \\
        \hline
        Deciding what is controllable & Rank of $\mathcal{C}$, $\mathcal{O}$ & Use \texttt{demo.py lqr --no-plot} to print ranks before attempting placement or LQR. \\
        \hline
        Optimal control + estimation & Discrete Riccati equations for LQR and Kalman & Run \texttt{EXPERIMENTS["lqr"]} and \texttt{EXPERIMENTS["kf\_lqr"]} to inspect gains, covariance logs, and innovation statistics. \\
        \hline
    \end{tabularx}
    \caption{The Module~\ref{chap:modern} workflow linking conceptual reasoning, algebra, and executable experiments.}
\end{table}

The presentation mirrors the step-by-step development advocated by \textcite{ljung2012control}, blending geometric insight with computational recipes.

\section{Why State-Space?}
\begin{beginnerguide}
    \item[\textbf{What?}] Explains the benefits of representing dynamics with internal state variables.
    \item[\textbf{Why?}] State-space handles multivariable systems, initial conditions, and modern control designs more naturally than transfer functions alone.
    \item[\textbf{When?}] Whenever systems have multiple inputs/outputs, need observer design, or involve optimal control.
    \item[\textbf{How?}] Define the state, write $\dot{x}=Ax+Bu$, $y=Cx+Du$, and relate it to familiar transfer-function descriptions.
    \item[\textbf{Example}] Convert a mass--spring--damper transfer function to controllable canonical form using \texttt{control.ss}.
\end{beginnerguide}
Transfer functions describe how inputs map to outputs by eliminating intermediate variables. State-space models keep those internal variables---the \emph{state}---explicit:
\begin{align}
    \dot{x} &= A x + B u,\\
    y &= C x + D u.
\end{align}
For beginners, it helps to answer three practical questions:
\begin{itemize}[leftmargin=*]
    \item \textbf{What is the state?} A minimal set of variables that completely describe the system at any instant. For a mass--spring--damper, $x = [\text{position};~\text{velocity}]$ is enough to predict future motion.
    \item \textbf{When should I use state-space?} Whenever you need to model multiple inputs/outputs, track initial conditions explicitly, or design modern controllers (state feedback, observers, optimal control). Classical transfer-function tools become cumbersome in these scenarios.
    \item \textbf{How does it relate to transfer functions?} If a system is linear and time invariant, you can move between the two representations. Eliminating the state from the equations produces the transfer function; factoring a transfer function back into first-order equations yields a state-space realisation.
\end{itemize}

\subsection{Connecting to Transfer Functions}
Consider a single-input single-output (SISO) transfer function
\[
    G(s) = \frac{b_{m}s^{m} + \dots + b_1 s + b_0}{s^{n} + a_{n-1} s^{n-1} + \dots + a_1 s + a_0},
\]
normalised so the leading denominator coefficient is $1$. One controllable canonical realisation is obtained by defining the state as the stacked derivatives of the output,
\[
    x = \begin{bmatrix} y \\ \dot{y} \\ \vdots \\ y^{(n-1)} \end{bmatrix},
\]
producing the matrices
\begin{align}
    A &= \begin{bmatrix}
        0 & 1 & 0 & \dots & 0 \\
        0 & 0 & 1 & \dots & 0 \\
        \vdots & \vdots & \vdots & \ddots & \vdots \\
        0 & 0 & 0 & \dots & 1 \\
        -a_0 & -a_1 & -a_2 & \dots & -a_{n-1}
    \end{bmatrix}, &
    B &= \begin{bmatrix} 0 \\ 0 \\ \vdots \\ 0 \\ 1 \end{bmatrix},\\
    C &= \begin{bmatrix} \beta_0 & \beta_1 & \dots & \beta_{n-1} \end{bmatrix}, &
    D &= d,
\end{align}
where the strictly proper case ($m \leq n-1$) sets $d = 0$ and $\beta_i = b_i$ (pad missing coefficients with zeros so the vector has length $n$). If the numerator order equals the denominator order, factor out the direct term $d = b_n$ and build the strictly proper remainder using the formula above. This construction guarantees that $G(s) = C (sI - A)^{-1} B + D$.

\subsection{Worked Example: Mass--Spring--Damper}
The plant $G(s) = \frac{1/m}{s^2 + \frac{c}{m} s + \frac{k}{m}}$ yields
\[
    A = \begin{bmatrix} 0 & 1 \\ -k/m & -c/m \end{bmatrix}, \quad
    B = \begin{bmatrix} 0 \\ 1/m \end{bmatrix}, \quad
    C = \begin{bmatrix} 1 & 0 \end{bmatrix}, \quad
    D = 0,
\]
matching the physics-based derivation in Chapter~\ref{chap:modeling}. The Python snippet
\begin{lstlisting}[language=python]
from control import tf, ss
G = tf([1], [1, 0.4, 4])      # m = 1, c = 0.4, k = 4
A, B, C, D = ss(G)
print(A)
print(B)
\end{lstlisting}
confirms the algebra numerically and reinforces the connection between the two viewpoints. Try modifying the transfer function (e.g., adding zeros) and watch how the canonical matrices change.

\section{Controllability and Observability}
\begin{beginnerguide}
    \item[\textbf{What?}] Criteria indicating whether inputs can move every state and outputs can reconstruct them.
    \item[\textbf{Why?}] Feedback and observers fail if the system is uncontrollable or unobservable; these checks prevent wasted design effort.
    \item[\textbf{When?}] Immediately after forming $(A,B,C,D)$ and before attempting pole placement or estimation.
    \item[\textbf{How?}] Compute the controllability matrix $\mathcal{C}$ and observability matrix $\mathcal{O}$ and inspect their ranks.
    \item[\textbf{Example}] Use \texttt{EXPERIMENTS["lqr"].runner()} to print the ranks for the discretised mass--spring--damper, confirming readiness for control design.
\end{beginnerguide}
The controllability matrix
\begin{equation}
    \mathcal{C} = [B\; AB\; A^2 B\; \dots\; A^{n-1} B]
\end{equation}
has full rank ($n$) if the state can be driven to any point in finite time. Observability relies on
\begin{equation}
    \mathcal{O} = \begin{bmatrix} C \\ C A \\ \vdots \\ C A^{n-1} \end{bmatrix},
\end{equation}
which must also have rank $n$ to reconstruct the state from outputs. Physically, uncontrollable states correspond to modes unaffected by actuators; unobservable states correspond to hidden modes in the measurements.

\subsection{Numerical Check}
Use the Python experiment
\begin{lstlisting}[language=bash]
python demo.py lqr --no-plot
\end{lstlisting}
which prints ranks of $\mathcal{C}$ and $\mathcal{O}$ for the discretized mass--spring--damper system, confirming controllability and observability.

\section{Pole Placement}
\begin{beginnerguide}
    \item[\textbf{What?}] A method for assigning closed-loop eigenvalues via state feedback.
    \item[\textbf{Why?}] It gives direct control over system dynamics when full-state feedback is available.
    \item[\textbf{When?}] After verifying controllability and before exploring optimal control alternatives like LQR.
    \item[\textbf{How?}] Use Ackermann’s formula or numerical solvers to compute $K$ such that $A-BK$ has desired eigenvalues.
    \item[\textbf{Example}] Call \texttt{scipy.signal.place\_poles} on a controllable canonical system to place poles at $[-2,-3]$ and compare simulated responses.
\end{beginnerguide}
Design a feedback law $u = -K x$ such that $A - B K$ has desired eigenvalues $\{\lambda_i\}$. Ackermann's formula for single-input systems yields
\begin{equation}
    K = e_n^T \mathcal{C}^{-1} \phi_d(A),
\end{equation}
where $e_n^T = [0\; 0\; \dots\; 1]$ and $\phi_d(A)$ denotes the desired characteristic polynomial evaluated at $A$. Practical workflow:
\begin{enumerate}[leftmargin=*]
    \item Choose desired poles based on transient requirements (e.g., match damping ratio and natural frequency).
    \item Use \texttt{scipy.signal.place\_poles} or control library to compute $K$ and verify eigenvalues numerically.
    \item Simulate closed-loop response to ensure no actuator saturation or excessive control effort occurs.
\end{enumerate}

\section{Linear Quadratic Regulation}
\begin{beginnerguide}
    \item[\textbf{What?}] An optimal control law balancing state deviations and control effort via the Riccati equation.
    \item[\textbf{Why?}] LQR offers systematic tuning knobs ($Q$, $R$) and guarantees closed-loop stability for controllable systems.
    \item[\textbf{When?}] After baseline pole placement when you want performance-effort trade-offs and robustness improvements.
    \item[\textbf{How?}] Solve the discrete or continuous algebraic Riccati equation, compute $K$, and optionally add a prefilter for tracking.
    \item[\textbf{Example}] Evaluate \texttt{EXPERIMENTS["lqr"].runner()} to see how different $Q$/$R$ choices alter transient response and control usage.
\end{beginnerguide}
LQR minimizes the infinite-horizon cost
\begin{equation}
    J = \sum_{k=0}^{\infty} (x_k^T Q x_k + u_k^T R u_k)
\end{equation}
subject to $x_{k+1} = A x_k + B u_k$. The solution involves:
\begin{enumerate}[leftmargin=*]
    \item Solving the discrete algebraic Riccati equation (DARE)
    \begin{equation}
        P = A^T P A - A^T P B (B^T P B + R)^{-1} B^T P A + Q.
    \end{equation}
    \item Computing $K = (B^T P B + R)^{-1} B^T P A$.
    \item Using $u_k = -K x_k$ to close the loop; the closed-loop eigenvalues equal those of $A - B K$ and lie inside the unit circle.
\end{enumerate}
Select $Q$ and $R$ to penalize state deviations and control effort. Increasing $Q$ entries tightens state regulation; increasing $R$ reduces control energy at the expense of slower response. To show where the Riccati equation originates, start from the finite-horizon cost
\[
J_N = x_N^T P_N x_N + \sum_{k=0}^{N-1} (x_k^T Q x_k + u_k^T R u_k)
\]
and minimise it by substituting $u_k = -K_k x_k$. Completing the squares yields the backward recursion
\begin{align*}
    P_k &= Q + A^T P_{k+1} A - A^T P_{k+1} B (R + B^T P_{k+1} B)^{-1} B^T P_{k+1} A,\\
    K_k &= (R + B^T P_{k+1} B)^{-1} B^T P_{k+1} A,
\end{align*}
with terminal condition $P_N = P_N^T \succeq 0$. As $N \to \infty$, the matrices converge to $P$ and $K$ that satisfy the steady-state DARE above. Add a short logging loop to the LQR experiment (or call the recursion interactively) to print each $P_k$ so beginners can watch the sequence settle numerically before trusting the steady-state solver.

\subsection{Servo Prefilter}
To track constant references without steady-state error, compute the prefilter gain
\begin{equation}
    N_{\text{bar}} = \frac{1}{C (I - A + B K)^{-1} B}
\end{equation}
for single-output systems. The helper \texttt{discrete\_nbar} automates this calculation. The control law becomes $u_k = -K x_k + N_{\text{bar}} r_k$.

\section{Kalman Filtering in Depth}
\begin{beginnerguide}
    \item[\textbf{What?}] Optimal state estimation under linear-Gaussian assumptions.
    \item[\textbf{Why?}] Accurate state estimates enable state-feedback controllers when not all states are measured and quantify uncertainty for diagnostics.
    \item[\textbf{When?}] After specifying noise covariances and before closing the loop with LQR or other controllers.
    \item[\textbf{How?}] Propagate predictions with $A$, $B$, update with measurements via the Kalman gain, and tune using innovation statistics.
    \item[\textbf{Example}] Run \texttt{EXPERIMENTS["kf\_lqr"].runner(return\_logs=True)} to collect innovations and verify tuning against the theoretical covariance.
\end{beginnerguide}
The Kalman filter is a probabilistic extension of the Luenberger observer that accounts for process and measurement uncertainty. We begin with the stochastic model
\begin{align}
    x_{k+1} &= A x_k + B u_k + w_k, & w_k &\sim \mathcal{N}(0, Q), \label{eq:kalman-process}\\
    y_k &= C x_k + v_k, & v_k &\sim \mathcal{N}(0, R), \label{eq:kalman-measure}
\end{align}
where $w_k$ and $v_k$ are mutually independent Gaussian sequences. The goal is to compute the conditional distribution of $x_k$ given all measurements $y_0,\dots,y_k$; under the linear-Gaussian assumptions this distribution remains Gaussian and is characterized completely by its mean $\hat{x}_{k|k}$ and covariance $P_{k|k}$.

\subsection{Projection Derivation}
Assume we know $\hat{x}_{k|k}$ and $P_{k|k}$ at time $k$. The one-step-ahead state prediction follows from taking expectations of \eqref{eq:kalman-process}:
\begin{align}
    \hat{x}_{k+1|k} &= A \hat{x}_{k|k} + B u_k,\\
    P_{k+1|k} &= A P_{k|k} A^\top + Q,
\end{align}
where $P_{k+1|k}$ is the a priori covariance. When the $(k+1)$-st measurement arrives, we correct the prediction by minimizing the mean-square estimation error:
\[
    \hat{x}_{k+1|k+1} = \hat{x}_{k+1|k} + K_{k+1} \bigl(y_{k+1} - C \hat{x}_{k+1|k}\bigr).
\]
Choosing $K_{k+1}$ to minimize $\operatorname{tr} P_{k+1|k+1}$ yields the familiar gain expression
\begin{equation}
    K_{k+1} = P_{k+1|k} C^\top \bigl(C P_{k+1|k} C^\top + R\bigr)^{-1}, \label{eq:kalman-gain}
\end{equation}
and the updated covariance
\begin{equation}
    P_{k+1|k+1} = (I - K_{k+1} C) P_{k+1|k}. \label{eq:kalman-cov-update}
\end{equation}
Equations \eqref{eq:kalman-process}--\eqref{eq:kalman-cov-update} constitute the discrete Kalman filter. Appendix~\ref{chap:math-appendix} contains the full derivation using conditional expectations and the orthogonality principle for learners interested in the rigorous proof.

\subsection{Innovation Interpretation}
Define the innovation (or residual)
\[
    r_{k+1} = y_{k+1} - C \hat{x}_{k+1|k}.
\]
It represents the new information brought by the latest measurement. Its covariance is
\[
    S_{k+1} = C P_{k+1|k} C^\top + R,
\]
which appears in \eqref{eq:kalman-gain}. Normalizing the innovations via
\[
    \nu_{k+1} = r_{k+1}^\top S_{k+1}^{-1} r_{k+1}
\]
provides a diagnostic statistic: for a correctly tuned filter, $\nu_{k+1}$ follows a $\chi^2$ distribution with degrees of freedom equal to the measurement dimension. Systematically large values indicate under-modelled process noise ($Q$ too small), while systematically small values imply the opposite.

\subsection{Steady-State Filter and Riccati Link}
If $(A, C)$ is detectable and $(A, Q^{1/2})$ is stabilizable, the Riccati difference equation converges to a positive-definite solution $P$ satisfying
\begin{equation}
    P = A P A^\top + Q - A P C^\top (C P C^\top + R)^{-1} C P A^\top. \label{eq:dares}
\end{equation}
Equating $P_{k+1|k} = P$ and $P_{k+1|k+1} = P$ in \eqref{eq:kalman-cov-update} recovers the discrete algebraic Riccati equation, underscoring the duality between optimal control and optimal estimation. The steady-state Kalman gain becomes
\[
    K = P C^\top (C P C^\top + R)^{-1},
\]
mirroring the LQR gain $K = (B^\top P B + R)^{-1} B^\top P A$ from the previous section.

\subsection{Practical Tuning Workflow}
\begin{enumerate}[leftmargin=*]
    \item \textbf{Scale the state.} Ensure each state component has comparable units; otherwise $Q$ requires manual rescaling.
    \item \textbf{Choose baseline covariances.} Start with $Q = q\,I$ where $q$ reflects typical unmodelled accelerations or torques, and set $R$ using the squared sensor standard deviation.
    \item \textbf{Run the filter and log innovations.} Use the built-in logging hook (see below) to capture $(t, r_k, S_k)$ pairs.
    \item \textbf{Diagnose.} If $\nu_k$ (normalized innovation squared) consistently exceeds the theoretical expectation, increase the corresponding entries in $Q$. If $\nu_k$ is too small, decrease $Q$ or increase $R$.
    \item \textbf{Validate closed-loop performance.} Re-run the coupled LQR--Kalman example and examine actuator usage and tracking error.
\end{enumerate}
The experiment harness exposes the innovations automatically when invoked with logging enabled:
\begin{lstlisting}[language=python]
from scripts.experiments import EXPERIMENTS
result = EXPERIMENTS["kf_lqr"].runner(return_logs=True)
innovation_log = result.metadata["innovation_log"]
t, innov, Sk = innovation_log.T  # unpack time series
\end{lstlisting}
Plot $|\text{innov}|$ against $\sqrt{S_k}$ to visualize whether the residuals stay within expected bounds.

\subsection{Case Study 1: Position-Only Sensing}
Consider the mass--spring--damper with position measurement only. Select
\[
    Q = \operatorname{diag}(10^{-4},\, 5\times 10^{-5}), \qquad R = 4\times 10^{-4}.
\]
Execute
\begin{lstlisting}[language=bash]
python demo.py kf_lqr --no-plot
\end{lstlisting}
and observe:
\begin{itemize}[leftmargin=*]
    \item convergence of the Kalman gain (velocity reconstruction emerges automatically);
    \item printed performance metrics showing zero steady-state error when combined with the LQR controller;
    \item innovation statistics that should match the chosen $Q$, $R$ if the tuning is consistent.
\end{itemize}
Increase $Q$ by a factor of ten and rerun the experiment. The estimator responds faster but injects more noise into the control effort---visible in \texttt{figures/modern\_lqr\_controls.png} after regenerating the tutorial plots. Conversely, reducing $Q$ slows the estimator yet yields smoother control.

\subsection{Case Study 2: Sensor Fusion with Bias}
Extend the model to include a sensor bias $b_k$ by augmenting the state:
\[
    x_k^\text{aug} = \begin{bmatrix} x_k \\ b_k \end{bmatrix}, \qquad
    b_{k+1} = b_k + \eta_k,\quad \eta_k \sim \mathcal{N}(0, \sigma_b^2).
\]
Modify \texttt{scripts/experiments.py} accordingly (see comments near the Kalman filter setup) and run the simulation with $\sigma_b = 5\times 10^{-5}$. The augmented filter learns the bias online, reducing steady-state error without sacrificing noise rejection. Compare the logged innovations before and after augmentation to verify that the residual statistics return to nominal levels.

\subsection{Case Study 3: Missing Data and Outliers}
Kalman filters handle occasional sensor dropouts by skipping the update step. Modify the controller loop to randomly skip updates with probability $p = 0.1$:
\begin{lstlisting}[language=python]
if np.random.rand() < 0.1:
    last_control = sf.step(kf.xhat, reference)  # prediction only
else:
    innovation = kf.update(measurement)
    last_control = sf.step(kf.xhat, reference)
\end{lstlisting}
Inspect the innovation log to ensure the covariance $P$ inflates during dropouts, reflecting increased uncertainty. Re-enable the update step when measurements resume and confirm the estimator smoothly re-converges. This exercise illustrates how the Kalman framework naturally extends to imperfect sensing environments, a common scenario in embedded control.


\section{Integrated Example: LQR with Kalman Filter}
\begin{beginnerguide}
    \item[\textbf{What?}] Combines optimal control and estimation into a full-state-feedback loop using the separation principle.
    \item[\textbf{Why?}] Demonstrates how controller and observer interact in practice, preparing you for implementation.
    \item[\textbf{When?}] After designing $K$ and Kalman gains separately, as a final integration step.
    \item[\textbf{How?}] Discretize the plant, compute LQR and Kalman gains, and simulate the coupled loop.
    \item[\textbf{Example}] Execute the provided experiment sequence to see tracking performance and estimator convergence for the mass--spring--damper.
\end{beginnerguide}
Run
\begin{lstlisting}[language=bash]
python demo.py kf_lqr
\end{lstlisting}
which performs the following steps:
\begin{enumerate}[leftmargin=*]
    \item Discretize the mass--spring--damper system with sampling time $T_s = 0.01$~s.
    \item Solve the DARE with $Q = \operatorname{diag}(150, 10)$, $R = 0.4$ to obtain $K$.
    \item Compute the prefilter $N_{\text{bar}}$ for step tracking.
    \item Run the Kalman filter with $Q_n = \operatorname{diag}(10^{-4}, 10^{-3})$ and $R_n = 10^{-4}$.
\end{enumerate}
Metrics such as steady-state error and overshoot are reported; compare them with the pure state-feedback case to appreciate estimator benefits.

\section{Tutorial: LQR and Kalman Practice}
\begin{beginnerguide}
    \item[\textbf{What?}] Guided scripting workflow that regenerates figures and metrics for LQR and Kalman designs.
    \item[\textbf{Why?}] Reinforces theory with reproducible plots and logs, making it easy to experiment with weighting matrices.
    \item[\textbf{When?}] After understanding the algorithms, as hands-on reinforcement or report preparation.
    \item[\textbf{How?}] Run \tutorialcmd[--output-dir figures]{modern}, inspect responses, and adjust $Q$, $R$, $Q_n$, $R_n$.
    \item[\textbf{Example}] Compare \texttt{figures/modern\_lqr\_responses.png} for “tracking” vs. “energy” weightings to see the effect of $Q$/$R$ scaling.
\end{beginnerguide}
\begin{enumerate}[leftmargin=*]
    \item Generate the modern-control figures with\\
    \tutorialcmd[--output-dir figures]{modern}.
    \item Use \texttt{figures/modern\_lqr\_responses.png} to compare "tracking" and "energy" weightings in the style suggested by \textcite{ljung2012control}; note how $Q$ and $R$ choices reshape the transient.
    \item Inspect \texttt{figures/modern\_lqr\_controls.png} to quantify the actuator effort required by each design and relate it to the penalty matrices.
    \item Study \texttt{figures/modern\_kalman\_tracking.png} to see the Kalman estimate converge toward the true trajectory despite measurement noise; experiment with $Q_n$, $R_n$ inside \texttt{tutorial\_modern} to observe innovation statistics.
\end{enumerate}
Exact commands and build automation tips are consolidated in Appendix~\ref{chap:tooling} for quick reference while you work through these exercises.
\IfFileExists{../figures/modern_lqr_responses.png}{
    \begin{figure}[h]
        \centering
        \includegraphics[width=0.75\linewidth]{../figures/modern_lqr_responses.png}
        \caption{LQR responses for different weightings produced by \tutorialcmd{modern}.}
    \end{figure}
}{
    \begin{figure}[h]
        \centering
        \fbox{\parbox{0.7\linewidth}{Run \tutorialcmd{modern} to create \texttt{figures/modern\_lqr\_responses.png}, then recompile.}}
        \caption{Placeholder for the LQR comparison plot.}
    \end{figure}
}
\IfFileExists{../figures/modern_lqr_controls.png}{
    \begin{figure}[h]
        \centering
        \includegraphics[width=0.75\linewidth]{../figures/modern_lqr_controls.png}
        \caption{Control effort comparison for alternative LQR designs.}
    \end{figure}
}{
    \begin{figure}[h]
        \centering
        \fbox{\parbox{0.7\linewidth}{Generate \texttt{figures/modern\_lqr\_controls.png} to compare control usage.}}
        \caption{Placeholder for the LQR control-effort plot.}
    \end{figure}
}
\IfFileExists{../figures/modern_kalman_tracking.png}{
    \begin{figure}[h]
        \centering
        \includegraphics[width=0.75\linewidth]{../figures/modern_kalman_tracking.png}
        \caption{Kalman-filtered tracking result generated by the modern tutorial script.}
    \end{figure}
}{
    \begin{figure}[h]
        \centering
        \fbox{\parbox{0.7\linewidth}{Generate \texttt{figures/modern\_kalman\_tracking.png} to visualize estimator performance.}}
        \caption{Placeholder for the Kalman tracking plot.}
    \end{figure}
}

\section{Design Trade-Offs}
\begin{beginnerguide}
    \item[\textbf{What?}] Discussion of how weighting choices, estimator tuning, and actuator limits interplay.
    \item[\textbf{Why?}] Understanding trade-offs prevents over-aggressive designs and highlights practical constraints.
    \item[\textbf{When?}] During iterative tuning and before deploying on hardware.
    \item[\textbf{How?}] Evaluate how changes in $Q$, $R$, $Q_n$, $R_n$ impact transient metrics, control effort, and innovation statistics.
    \item[\textbf{Example}] Increase $Q$ entries in the tutorial script, observe reduced settling time but higher control effort, and decide acceptable compromises.
\end{beginnerguide}
\begin{itemize}[leftmargin=*]
    \item \textbf{Weight selection}: Large $Q$ entries can yield aggressive state regulation but may demand high control effort, risking actuator saturation. Balance with $R$.
    \item \textbf{Estimator tuning}: Underestimating measurement noise $R$ leads to noisy control inputs; overestimating $Q$ yields sluggish response. Use innovation analysis to refine values.
    \item \textbf{Reference tracking}: For non-zero steady-state goals, augment the system with integral action or design feedforward prefilters.
\end{itemize}

\begin{takeaway}
\begin{itemize}[leftmargin=*]
    \item Controllability and observability form the foundation of modern control; verify them before attempting advanced designs.
    \item LQR provides a principled way to balance performance and effort, while prefilters ensure accurate tracking.
    \item Kalman filtering enables high-performance control despite sensor noise, completing the state-estimation feedback loop.
\end{itemize}
\end{takeaway}

\section{Module Review}
\begin{itemize}[leftmargin=*]
    \item \textbf{Key tests}: controllability matrix $\mathcal{C}$ and observability matrix $\mathcal{O}$ must be full rank before designing feedback or estimators.
    \item \textbf{LQR summary}: solve the DARE for $P$, compute $K = (B^T P B + R)^{-1} B^T P A$, and use $N_{\text{bar}}$ for steady-state accuracy; interpret $Q$, $R$ weightings with the tutorial overlay in \texttt{figures/modern\_lqr\_responses.png}.
    \item \textbf{Kalman filter}: tune $Q_n$, $R_n$ using innovation statistics; verify estimator response with Figure~\ref{fig:adaptive-tracking} and \texttt{figures/modern\_kalman\_tracking.png}.
    \item \textbf{Practice}: regenerate \texttt{figures/modern\_lqr\_responses.png} and modify $Q$, $R$ in \texttt{tutorial\_modern} to see trade-offs described in \textcite{ljung2012control}.
    \item \textbf{Reading}: \textcite{ljung2012control} Ch.~9--11 for identification-enabled state-space design and \textcite{astrom2006advanced} for practical loop-shaping comparisons.
\end{itemize}

\section{Keyword Review}
\begin{vocabulary}
\begin{itemize}[leftmargin=*]
    \item \textbf{State}: minimal set of internal variables required to determine future evolution with the input.
    \item \textbf{Controllability}: ability to reach any state via inputs; captured by the rank of $\mathcal{C}$.
    \item \textbf{Observability}: ability to reconstruct the state from outputs; captured by the rank of $\mathcal{O}$.
    \item \textbf{LQR}: optimal state-feedback law obtained by solving a Riccati equation to balance error and control effort.
    \item \textbf{Kalman gain}: $K = P C^\top (C P C^\top + R)^{-1}$, weighting how much the estimate trusts new measurements.
    \item \textbf{Separation principle}: controller and observer can be designed independently when the system is controllable and observable.
\end{itemize}
\end{vocabulary}
